\documentclass[10pt,a4paper]{amsart}
\usepackage[margin=19mm]{geometry}
\usepackage{amsmath,amssymb,mathtools}
\usepackage[scr=rsfs,cal=euler]{mathalfa}
\usepackage{xfrac}
\usepackage[unicode,hidelinks,bookmarks=false]{hyperref}
\newcommand{\ie}{i.e.\ }
\newcommand{\cf}{cf.\ }
\newcommand{\resp}{respectively\ }
\newcommand{\Subsection}{\subsection}
\providecommand{\nobreakdash}{\nobreak\hbox{-}}
\setlength{\emergencystretch}{2em}
\multlinegap0pt
\usepackage{amssymb}
\usepackage{bm}
\usepackage{dsfont}
\usepackage{xcolor}
\usepackage[normalem]{ulem}
\usepackage{tikz}
\newtheorem{lemma}{Lemma}
\newtheorem{theorem}{Theorem}
\usepackage{cleveref}
\crefname{lemma}{Lemma}{Lemmas}
\crefname{theorem}{Theorem}{Theorems}
\title{Completely-positive non-signalling non-Markovian dynamics}
\author{Serhii Kryhin, Vivishek Sudhir}
\date{Source: arXiv:2605.03197v2 (CC BY 4.0)}
\begin{document}
\maketitle

\noindent\emph{Source and licence.} Completely-positive non-signalling non-Markovian dynamics. Version arXiv:2605.03197v2 (CC BY 4.0), distributed by arXiv under the Creative Commons Attribution 4.0 International licence, http://creativecommons.org/licenses/by/4.0/, as recorded in the arXiv licence metadata for this exact version and shown on the abstract page https://arxiv.org/abs/2605.03197v2. Full licence notice: \texttt{licenses/arxiv-2605.03197-CC-BY-4.0.txt}.\par\smallskip

\noindent\textit{Editorial scope.} Counted unit: the article's theorem thm:ConeExtension (source0.tex lines 706--712). The article's complete original proof of it is delivered with it (source0.tex lines 714--764, one \texttt{proof} environment of 51 lines). No mathematics is written, repaired or re-derived: the statement and the proof are the article's own lines, in the article's own notation, and no retained line is substituted or re-flowed.  Retained uncounted as the source-local context the proof needs: lines 695--697.  Nothing was omitted from any retained span, so the package carries no omission record.  No retained line contains a citation, so the wrapper carries no bibliography and no bibliography entry.  No retained line contains a figure environment, an image-inclusion command or drawing code, so no image is delivered and the package declares no asset dependency.  Editorial formatting only, in addition: an AMS \texttt{amsart} preamble that restates the article's own declarations and macros used by the retained lines, namely the environments \texttt{lemma}, \texttt{theorem} on separate counters, together with the standard packages those lines need.  The retained environments print in this document as Lemma 1, Theorem 1 in order of appearance, whereas the article prints the same items with the numbering of its own full document.  The complete native source of the article as distributed in the arXiv e-print for this exact version is bundled unchanged as \texttt{sources/199775/source0.tex}.
\begin{lemma}\label{lem:UniqueDecomp}
	Every nonzero $x\in C$ can be uniquely expressed as $x = \alpha b$ for some $\alpha > 0$ and $b \in B_f(C)$.
\end{lemma}

\begin{theorem}\label{thm:ConeExtension}
Let $C$ be a convex cone with a base $B_f(C)$, and let $h: B_f(C) \rightarrow B_f(C)$ be a bounded affine
map. Then there exists a unique bounded linear map $\tilde h$ on the cone such that 
(i) $\tilde h(\alpha x) = \alpha \tilde{h}(x)$ for all $x\in C$ and $\alpha > 0$; 
(ii) $\tilde h\vert_{B_f(C)} = h$; and 
(iii) $f(\tilde h(x)) = f(x)$ for all $x \in C$. 
\end{theorem}

\begin{proof}
Define $\tilde h: C \to C$ by 
\begin{equation*}
	\tilde h(x) = f(x)\, h\!\left(\frac{x}{f(x)}\right).
\end{equation*}
This is well-defined: $f(x) > 0$ implies $x/f(x) \in B_f(C)$, so $h(x/f(x)) \in B_f(C) \subset C$.

\noindent The boundedness of $\tilde h$ follows from $\| \tilde h (x) \| = |f(x)| \| h(x/f(x))\| \leq C_h \| x\|$, where $C_h$ is the norm of $h$.

\noindent The conditions (i), (ii), and (iii) are easy to verify:

(i) For $\alpha > 0$: $\tilde h(\alpha x) = f(\alpha x)\, h\!\left(\frac{\alpha x}{f(\alpha x)}\right) = \alpha f(x)\, h\!\left(\frac{x}{f(x)}\right) = \alpha\, \tilde h(x)$,
using the linearity of $f$.

(ii) For $b \in B_f(C)$: $\tilde h(b) = 1 \cdot h(b) = h(b)$.

(iii) $f(\tilde h(x)) = f\!\left(f(x)\, h\!\left(\frac{x}{f(x)}\right)\right) = f(x)\, f\!\left(h\!\left(\frac{x}{f(x)}\right)\right) = f(x)$,
since $h(x/f(x)) \in B_f(C)$ implies $f(h(x/f(x))) = 1$.

\noindent \emph{Linearity.} We first show $\tilde h(x+y) = \tilde h(x) + \tilde h(y)$ for $x,y \in C$.
Set $\alpha = f(x)$, $\beta = f(y)$, so that $x/\alpha,\, y/\beta \in B_f(C)$.
Since $C$ is a convex cone, $x+y \in C$, and by linearity of $f$, $f(x+y) = \alpha + \beta$; thus 
by the definition of $\tilde h$:
\begin{equation*}
\tilde h(x+y) = (\alpha + \beta)\, h\!\left(\frac{x+y}{\alpha + \beta}\right).
\end{equation*}
The last factor can be simplified by writing
\begin{equation*}
	\frac{x+y}{\alpha + \beta} = \underbrace{\frac{\alpha}{\alpha+\beta}}_{\lambda}\frac{x}{\alpha} 
	+ \underbrace{\frac{\beta}{\alpha+\beta}}_{1-\lambda}\frac{y}{\beta}
\end{equation*}
and using the affineness of $h$:
\begin{equation*}
\begin{split}
	\tilde h(x+y) %&= (\alpha+\beta)\, h\!\left(\frac{x+y}{\alpha+\beta}\right) \\
	= (\alpha+\beta)\left[\frac{\alpha}{\alpha+\beta}\, h\!\left(\frac{x}{\alpha}\right) 
	+ \frac{\beta}{\alpha+\beta}\, h\!\left(\frac{y}{\beta}\right)\right]
	= \alpha\, h\!\left(\frac{x}{\alpha}\right) + \beta\, h\!\left(\frac{y}{\beta}\right) 
	= \tilde h(x) + \tilde h(y).
\end{split}
\end{equation*}
Combined with (i), this gives $\tilde h(\alpha x + \beta y) = \alpha\,\tilde h(x) + \beta\,\tilde h(y)$ for all $\alpha,\beta > 0$, so $\tilde h$ is linear on $C$.

\noindent \emph{Uniqueness} follows from \cref{lem:UniqueDecomp} and (i), (ii):
let $h': C \rightarrow C$ be any linear map satisfying (i), (ii), then
\begin{equation*}
	h'(x) = h'\!\left(f(x)\cdot\frac{x}{f(x)}\right) 
	\overset{(\mathrm{i})}{=} f(x)\, h'\!\left(\frac{x}{f(x)}\right) 
	\overset{(\mathrm{ii})}{=} f(x)\, h\!\left(\frac{x}{f(x)}\right) = \tilde h(x).
\end{equation*}
\end{proof}

\end{document}
